\scanpage{131}
Supposons d'abord $p<1$, alors on a, d'après
une formule de majoration classique pour la
somme des puissances $p$\textsuperscript{ièmes} des inverses des
modules des zéros d'une fonction entière
\[
f(z)=\det(1-zA)\qquad(0<p<1)
\]
\[
s_p(A)\leq p^2\int_0^\infty\frac{\log M_A(r)}{r^{p+1}}\,dr\tag*{\deleted{(10)}}
\]
or $M_A(\unsure{v|r})\leq M_{|A|}(\unsure{v|r})$ (\unsure{\ldots\ --\ --}) d'où
\[
s_p(A)\leq p^2\int_0^\infty\frac{\log M_{|A|}(\unsure{v|r})}{{\unsure{v|r}}^{\unsure{p-1|p+1}}}\,dr
\]
\,; appliquant alors la formule
(9), on obtient donc la majoration
\[
s_p(A)\leq\frac{p\pi}{\sin p\pi}S_p(A)\tag*{(10)}
\]
Le facteur $\frac{p\pi}{\sin p\pi}$ est $>1$, mais tend vers 1
pour $p\to0$. Soit alors $0<p\unsure{\leq|<}+\infty$ quelconque,
et soit $n$ un entier $>0$ tel que $\frac{p}{n}<1$.
On a évidemment $s_p(A)=s_{p/n}(A^n)$, d'où, en
vertu de (10) :
\[
s_p(A)\leq\frac{\frac{p}{n}\pi}{\sin\frac{p}{n}\pi}S_{p/n}(A^n)
\]
Or, on a \deleted{\unsure{$S_{p/n}(A^n)\leq S_p(A)$}}
\[
\|A^n\|_{p/n}\leq\underbrace{\|A\|_p\ldots\|A\|_p}_{n\ \text{facteurs}},
\]
comme il résulte aussitôt de la formule (6), qui sera.
démontrée plus bas (alors il est immédiat
que cette formule se généralise au produit d'un
nb fini quelconque d'opérateurs), d'où
\[
S_{p/n}(A^{\unsure{{-}|n}})=\left(\|A^n\|_{p/n}\right)^{p/n}
\leq\|A\|_p^p=S_p(A).
\]
